%===================================================================
\section{Generalized Defect-Cluster Population Dynamics}
\label{sec:general_defect_cluster_dynamics}
%===================================================================

This section develops a generalized master-equation formulation for the
population dynamics of defect clusters in irradiated crystalline solids. The
formulation is intended to describe the production of vacancies,
self-interstitial atoms, and gas atoms by irradiation; the subsequent formation
and evolution of interstitial and vacancy-cluster populations; the trapping and
release of gas atoms by vacancy clusters; and cross-species collisions between
interstitial and vacancy clusters. The formulation is written first in an
abstract reaction-network form, then in a code-generation-ready form, and
finally specialized to the vacancy--interstitial--helium system.

The resolved defect-cluster system contains two primary species classes:
\begin{equation}
    \mathcal{C}_I
    =
    \text{interstitial clusters},
    \qquad
    \mathcal{C}_V
    =
    \text{vacancy clusters/cavities}.
\end{equation}
Vacancy clusters are also referred to as \emph{cavities}. A cavity may be an
empty void or a gas-bearing bubble. Gas atoms that are not trapped in cavities
are represented by a matrix gas reservoir. Gas atoms do not form a third
resolved cluster class in the present formulation, and gas--gas reactions are
not allowed.

%===================================================================
\subsection{Primitive Species and Gas Reservoir}
\label{sec:primitive_species_and_gas_reservoir}
%===================================================================

Let the primitive irradiation-produced species be
\begin{equation}
    \mathcal{E}=\{I,V,\mathbf{g}\},
\end{equation}
where $I$ denotes a self-interstitial atom, $V$ denotes a vacancy, and
$\mathbf{g}$ denotes an ordered list of gas species. The gas list is written as
\begin{equation}
    \mathbf{g}
    =
    (g_1,g_2,\ldots,g_{N_g}),
    \label{eq:gas_species_list}
\end{equation}
where $N_g$ is the number of gas species included in the model. Depending on
the physical application,
\begin{equation}
    \mathbf{g}=\varnothing,
    \qquad
    \mathbf{g}=(\mathrm{He}),
    \qquad
    \mathbf{g}=(\mathrm{H},\mathrm{T}),
    \qquad
    \mathbf{g}=(\mathrm{Xe}),
\end{equation}
or a larger gas-species list may be used.

Gas atoms in the matrix are represented by reservoir concentrations
\begin{equation}
    c^g_\mu(t),
    \qquad
    \mu=1,\ldots,N_g,
    \label{eq:gas_reservoir_concentration}
\end{equation}
where $\mu$ labels the gas species $g_\mu$. These variables represent
monomeric gas atoms in the matrix reservoir. They are coupled to vacancy
clusters through trapping, emission, re-solution, and gas release from
eliminated cavities. The model does not include reactions of the type
\begin{equation}
    g_\mu+g_\nu\longrightarrow \text{gas cluster},
\end{equation}
so gas-only clusters are excluded.

%===================================================================
\subsection{Species Classes, Defect Populations, and State Variables}
\label{sec:species_classes_populations_state_variables}
%===================================================================

The formulation distinguishes between a \emph{species class} and a
\emph{defect population}. A species class identifies the broad physical type of
defect cluster. A defect population identifies a subcollection of clusters
within a species class that share the same mobility law, orientation state,
microstructural environment, or sink coupling.

The interstitial-cluster species class is denoted by $\mathcal{C}_I$. Its
defect populations are denoted by
\begin{equation}
    \mathcal{P}_I^\alpha \subset \mathcal{C}_I,
    \qquad
    \alpha=1,\ldots,N_P^I .
\end{equation}
The label $\alpha$ is a population label, not a cluster-size index. For
example, different interstitial-cluster populations may represent mobile SIA
clusters, sessile SIA loops, loops aligned with a principal stress direction,
or loops not aligned with that direction.

The vacancy-cluster species class is denoted by $\mathcal{C}_V$. Its defect
populations are denoted by
\begin{equation}
    \mathcal{P}_V^\beta \subset \mathcal{C}_V,
    \qquad
    \beta=1,\ldots,N_P^V .
\end{equation}
The label $\beta$ is also a population label. For example, different
vacancy-cluster populations may represent bulk cavities, dislocation-attached
cavities, grain-boundary cavities, or surface-attached cavities.

An interstitial-cluster state is written as
\begin{equation}
    I_{n,\alpha},
    \qquad
    n=1,\ldots,N_I,
    \qquad
    \alpha=1,\ldots,N_P^I,
    \label{eq:I_state}
\end{equation}
where $n$ is the number of self-interstitial atoms and $\alpha$ identifies the
interstitial-cluster population. Its concentration is
\begin{equation}
    c^I_{n,\alpha}(t).
    \label{eq:I_concentration}
\end{equation}

A vacancy-cluster, or cavity, state is written as
\begin{equation}
    V_{m,\boldsymbol{\ell},\beta},
    \qquad
    m=1,\ldots,N_V,
    \qquad
    \boldsymbol{\ell}\in\mathcal{G}_m,
    \qquad
    \beta=1,\ldots,N_P^V .
    \label{eq:V_state}
\end{equation}
Here $m$ is the number of vacancies in the cavity, $\beta$ identifies the
vacancy-cluster population, and
\begin{equation}
    \boldsymbol{\ell}
    =
    (\ell_1,\ell_2,\ldots,\ell_{N_g})
    \in \mathbb{Z}_{\geq 0}^{N_g}
    \label{eq:gas_occupancy_vector}
\end{equation}
is the gas inventory trapped inside the cavity. The component $\ell_\mu$ is
the number of atoms of gas species $g_\mu$ trapped inside the cavity. The set
$\mathcal{G}_m$ denotes the admissible gas occupancies for a cavity containing
$m$ vacancies. Its concentration is
\begin{equation}
    c^V_{m,\boldsymbol{\ell},\beta}(t).
    \label{eq:V_concentration}
\end{equation}

The notation $V_{m,\boldsymbol{\ell},\beta}$ denotes a vacancy-cluster, or cavity, state. Conventional chemical notation may be recovered as a
special case. For example, if $\mathbf{g}=(\mathrm{He})$, then
\begin{equation}
    V_{m,(\ell),\beta}
    \quad\text{denotes}\quad
    V_m\mathrm{He}_{\ell}
    \quad\text{in population }\beta .
\end{equation}
If $\mathbf{g}=(\mathrm{He},\mathrm{T})$, then
\begin{equation}
    V_{m,(\ell,p),\beta}
    \quad\text{denotes}\quad
    V_m\mathrm{He}_{\ell}\mathrm{T}_{p}
    \quad\text{in population }\beta .
\end{equation}
A gas-free cavity is obtained by setting
\begin{equation}
    \boldsymbol{\ell}=\mathbf{0},
    \qquad
    V_{m,\mathbf{0},\beta}
    \quad\text{denotes a pure vacancy cluster }V_m
    \quad\text{in population }\beta .
\end{equation}

The full deterministic state vector is therefore
\begin{equation}
    \mathbf{c}(t)
    =
    \left[
    \{c^I_{n,\alpha}\}_{n,\alpha},
    \{c^V_{m,\boldsymbol{\ell},\beta}\}_{m,\boldsymbol{\ell},\beta},
    \{c^g_\mu\}_{\mu=1}^{N_g}
    \right]^\top .
    \label{eq:general_state_vector}
\end{equation}
The first block contains interstitial-cluster populations, the second block
contains vacancy-cluster/cavity populations, and the third block contains
matrix gas-reservoir variables.

%===================================================================
\subsection{Signed Lattice-Defect Content}
\label{sec:signed_lattice_defect_content}
%===================================================================

Vacancies and interstitials represent opposite deviations from the perfect
lattice. This motivates the use of a signed lattice-defect content. For the
resolved states,
\begin{equation}
    q(I_{n,\alpha})=+n,
    \qquad
    q(V_{m,\boldsymbol{\ell},\beta})=-m,
    \qquad
    q(g_\mu)=0 .
    \label{eq:signed_defect_content}
\end{equation}
The population labels $\alpha$ and $\beta$ do not affect the signed defect
content; they affect mobility, sink coupling, reaction rates, and admissible
population-transfer pathways.

Cross-species collisions between interstitial and vacancy clusters are
therefore described by signed-defect cancellation. For a collision between an
interstitial cluster $I_{n,\alpha}$ and a vacancy cluster
$V_{m,\boldsymbol{\ell},\beta}$, the outgoing signed defect content is
\begin{equation}
    q_{\mathrm{out}} = n-m .
\end{equation}
Thus the lattice-defect part of the product is an interstitial cluster if
$n>m$, no lattice defect if $n=m$, and a vacancy cluster if $m>n$.

The generic signed-collision rule is
\begin{equation}
    I_{n,\alpha}
    +
    V_{m,\boldsymbol{\ell},\beta}
    \longrightarrow
    \begin{cases}
    I_{n-m,\alpha'} + \mathcal{R}_{\boldsymbol{\ell}}, & n>m, \\[4pt]
    \mathcal{R}_{\boldsymbol{\ell}}, & n=m, \\[4pt]
    V_{m-n,\boldsymbol{\ell},\beta'}, & m>n .
    \end{cases}
    \label{eq:signed_collision_rule}
\end{equation}
Here $\alpha'$ and $\beta'$ denote the product populations, determined by the
specific reaction rule. In the simplest single-population model,
$\alpha'=\alpha$ and $\beta'=\beta$ may be suppressed. The operator
$\mathcal{R}_{\boldsymbol{\ell}}$ denotes release of the gas inventory from an
eliminated cavity to the matrix gas reservoir. In the default reservoir model,
\begin{equation}
    \mathcal{R}_{\boldsymbol{\ell}}
    =
    \sum_{\mu=1}^{N_g} \ell_\mu g_\mu .
    \label{eq:gas_release_operator}
\end{equation}
If $m>n$, the cavity survives and retains its gas inventory
$\boldsymbol{\ell}$. If $n\geq m$, the supporting vacancy cluster is
eliminated and its gas content is returned to the gas reservoir.

%===================================================================
\subsection{Reaction-Network Representation}
\label{sec:reaction_network_representation}
%===================================================================

The master equation is assembled from elementary reactions. Each reaction is
assigned an index $r$. The index $r$ labels one specific physical process, for
example interstitial--interstitial coalescence, vacancy--vacancy coalescence,
gas trapping, gas emission, trap mutation, radiation re-solution,
population transfer, sink absorption, or an interstitial--vacancy collision.

Let $X_a$ denote a generic state in the model. The state index $a$ may refer
to an interstitial cluster, a vacancy cluster, or a gas-reservoir species:
\begin{equation}
    X_a
    \in
    \left\{
    I_{n,\alpha},\;
    V_{m,\boldsymbol{\ell},\beta},\;
    g_\mu
    \right\}.
    \label{eq:generic_state_index}
\end{equation}
Thus $a$ is not a physical size by itself. It is a compact label for one state
variable in the ODE system.

For a given reaction $r$, let
\begin{equation}
    \nu^-_{ar}
\end{equation}
be the number of objects of state $X_a$ consumed by the reaction, and let
\begin{equation}
    \nu^+_{ar}
\end{equation}
be the number of objects of state $X_a$ produced by the reaction. The net
change of state $X_a$ caused by one occurrence of reaction $r$ is
\begin{equation}
    S_{ar}
    =
    \nu^+_{ar}-\nu^-_{ar} .
    \label{eq:stoich_increment}
\end{equation}
The quantity $S_{ar}$ is the stoichiometric coefficient of state $a$ in
reaction $r$.

A general reaction can therefore be written compactly as
\begin{equation}
    \sum_a \nu^-_{ar} X_a
    \longrightarrow
    \sum_a \nu^+_{ar} X_a .
    \label{eq:general_reaction_network_form}
\end{equation}
This expression means that reaction $r$ consumes all states with nonzero
$\nu^-_{ar}$ and produces all states with nonzero $\nu^+_{ar}$.

For example, the interstitial-cluster coalescence reaction
\begin{equation}
    I_{n,\alpha}+I_{n',\alpha'}
    \longrightarrow
    I_{n+n',\alpha''}
    \label{eq:example_II_stoich}
\end{equation}
has
\begin{equation}
    S_{I_{n,\alpha},r}=-1,
    \qquad
    S_{I_{n',\alpha'},r}=-1,
    \qquad
    S_{I_{n+n',\alpha''},r}=+1,
\end{equation}
with all other entries of $S_{ar}$ equal to zero. Similarly, gas trapping by a
cavity,
\begin{equation}
    g_\mu + V_{m,\boldsymbol{\ell},\beta}
    \longrightarrow
    V_{m,\boldsymbol{\ell}+\mathbf{e}_\mu,\beta'},
    \label{eq:example_gas_trapping_stoich}
\end{equation}
has
\begin{equation}
    S_{g_\mu,r}=-1,
    \qquad
    S_{V_{m,\boldsymbol{\ell},\beta},r}=-1,
    \qquad
    S_{V_{m,\boldsymbol{\ell}+\mathbf{e}_\mu,\beta'},r}=+1.
\end{equation}
Here $\mathbf{e}_\mu$ is the unit vector in gas-species space.

%===================================================================
\subsection{Deterministic Master Equation}
\label{sec:deterministic_master_equation}
%===================================================================

The deterministic concentration of state $X_a$ is denoted by $c_a(t)$. In
terms of the explicit variables,
\begin{equation}
    c_a
    \in
    \left\{
    c^I_{n,\alpha},\;
    c^V_{m,\boldsymbol{\ell},\beta},\;
    c^g_\mu
    \right\}.
\end{equation}
The deterministic mean-field master equation is
\begin{equation}
    \boxed{
    \frac{dc_a}{dt}
    =
    G_a
    +
    \sum_r S_{ar}J_r(\mathbf{c})
    -
    D_a c_a .
    }
    \label{eq:abstract_master_equation}
\end{equation}
Here $G_a$ is the direct production rate of state $a$, $J_r(\mathbf{c})$ is
the flux of reaction $r$, and $D_a c_a$ is a first-order loss to unresolved
fixed sinks. Equation~\eqref{eq:abstract_master_equation} states that the
concentration of state $a$ changes by summing the contributions of all
reactions that produce or consume that state.

The reaction flux $J_r$ depends on the reaction type. For a unary transition,
\begin{equation}
    X_b \longrightarrow \text{products},
    \qquad
    J_r(\mathbf{c}) = \Gamma_r c_b,
    \label{eq:unary_flux}
\end{equation}
where $\Gamma_r$ is a first-order rate coefficient. For a binary collision,
\begin{equation}
    X_b+X_c \longrightarrow \text{products},
    \qquad
    J_r(\mathbf{c}) = K^{(r)}_{bc} c_b c_c,
    \label{eq:binary_flux}
\end{equation}
where $K^{(r)}_{bc}$ is the appropriate collision or capture rate. For a direct
source process,
\begin{equation}
    \emptyset \longrightarrow X_a,
    \qquad
    J_r=G_r .
    \label{eq:source_flux}
\end{equation}

%===================================================================
\subsection{Elementary Process Classes}
\label{sec:elementary_process_classes}
%===================================================================

The generalized model is built from the following process classes.

\subsubsection*{Irradiation production.}

Displacement cascades and gas-production processes introduce source terms
\begin{equation}
    G_a
    =
    G_a^{\mathrm{FP}}
    +
    G_a^{\mathrm{cascade}}
    +
    G_a^{\mathrm{gas}}
    +
    G_a^{\mathrm{trans}},
    \label{eq:general_source_decomposition}
\end{equation}
where the terms denote Frenkel-pair production, clustered cascade production,
direct gas production, and transmutation or implantation sources. Conservative
Frenkel-pair production satisfies
\begin{equation}
    \sum_{n,\alpha} n G^I_{n,\alpha}
    =
    \sum_{m,\boldsymbol{\ell},\beta} m G^V_{m,\boldsymbol{\ell},\beta} .
    \label{eq:cascade_pair_conservation}
\end{equation}

\subsubsection*{Interstitial--interstitial coalescence.}

Interstitial clusters may coalesce according to
\begin{equation}
    I_{n,\alpha}+I_{n',\alpha'}
    \longrightarrow
    I_{n+n',\alpha''},
    \label{eq:II_coalescence}
\end{equation}
with flux
\begin{equation}
    J^{II}_{n\alpha,n'\alpha'}
    =
    K^{II}_{n\alpha,n'\alpha'}
    c^I_{n,\alpha}c^I_{n',\alpha'} .
\end{equation}
The product population $\alpha''$ is specified by the physical population rule.
For a single interstitial population, the population indices are suppressed.

\subsubsection*{Vacancy--vacancy coalescence.}

Vacancy clusters may coalesce according to
\begin{equation}
    V_{m,\boldsymbol{\ell},\beta}
    +
    V_{m',\boldsymbol{\ell}',\beta'}
    \longrightarrow
    V_{m+m',\boldsymbol{\ell}+\boldsymbol{\ell}',\beta''} .
    \label{eq:VV_coalescence}
\end{equation}
The vacancy numbers add, and the gas occupancy vectors add componentwise. The
product population $\beta''$ is determined by the population rule for the
coalescence event.

\subsubsection*{Interstitial--vacancy signed cancellation.}

Cross-species collisions between interstitial and vacancy clusters obey
Eq.~\eqref{eq:signed_collision_rule}. These events conserve signed
lattice-defect content and conserve each gas species after the matrix reservoir
is included.

\subsubsection*{Gas trapping by cavities.}

A gas monomer of species $\mu$ may be trapped by a vacancy cluster:
\begin{equation}
    g_\mu
    +
    V_{m,\boldsymbol{\ell},\beta}
    \longrightarrow
    V_{m,\boldsymbol{\ell}+\mathbf{e}_\mu,\beta'} .
    \label{eq:gas_trapping}
\end{equation}
The corresponding flux is
\begin{equation}
    J^{\mathrm{trap}}_{m\boldsymbol{\ell}\beta,\mu}
    =
    K^{\mathrm{trap}}_{m\boldsymbol{\ell}\beta,\mu}
    c^g_\mu c^V_{m,\boldsymbol{\ell},\beta} .
\end{equation}

\subsubsection*{Gas emission and radiation re-solution.}

Gas may be thermally emitted or athermally re-solved from a cavity:
\begin{equation}
    V_{m,\boldsymbol{\ell},\beta}
    \longrightarrow
    V_{m,\boldsymbol{\ell}-\mathbf{e}_\mu,\beta'}
    +
    g_\mu,
    \qquad
    \ell_\mu\geq 1 .
    \label{eq:gas_emission_general}
\end{equation}
The corresponding flux may include both thermal and athermal contributions:
\begin{equation}
    J^{\mathrm{rel}}_{m\boldsymbol{\ell}\beta,\mu}
    =
    \left[
    \varepsilon_\mu(m,\boldsymbol{\ell},\beta)
    +
    \Gamma^{\mathrm{res}}_\mu(m,\boldsymbol{\ell},\beta)
    \right]
    c^V_{m,\boldsymbol{\ell},\beta} .
    \label{eq:gas_release_flux}
\end{equation}

\subsubsection*{Trap mutation.}

An over-pressurized gas-bearing cavity may emit an SIA while increasing its
vacancy content:
\begin{equation}
    V_{m,\boldsymbol{\ell},\beta}
    \longrightarrow
    V_{m+1,\boldsymbol{\ell},\beta'}
    +
    I_{1,\alpha'} .
    \label{eq:trap_mutation_general}
\end{equation}
This reaction conserves signed lattice-defect content because
\begin{equation}
    -m = -(m+1)+1 .
\end{equation}

\subsubsection*{Thermal vacancy and interstitial emission.}

Vacancy emission from a cavity is represented by
\begin{equation}
    V_{m,\boldsymbol{\ell},\beta}
    \longrightarrow
    V_{m-1,\boldsymbol{\ell},\beta'}
    +
    V_{1,\mathbf{0},\beta''},
    \qquad
    m\geq 2 .
    \label{eq:vacancy_emission_general}
\end{equation}
SIA emission from an interstitial cluster is represented by
\begin{equation}
    I_{n,\alpha}
    \longrightarrow
    I_{n-1,\alpha'}+I_{1,\alpha''},
    \qquad
    n\geq 2 .
    \label{eq:SIA_emission_general}
\end{equation}
Both processes conserve signed lattice-defect content.

\subsubsection*{Population transfer.}

A cluster may move between defect populations without changing its size or gas
inventory. For interstitial clusters,
\begin{equation}
    I_{n,\alpha}
    \longrightarrow
    I_{n,\alpha'},
    \label{eq:I_population_transfer}
\end{equation}
and for vacancy clusters,
\begin{equation}
    V_{m,\boldsymbol{\ell},\beta}
    \longrightarrow
    V_{m,\boldsymbol{\ell},\beta'} .
    \label{eq:V_population_transfer}
\end{equation}
These transitions may represent changes in mobility, stress orientation,
attachment to a dislocation, capture by a grain boundary, or transfer between
bulk and microstructural reservoirs. They conserve signed defect content and
gas inventory.

\subsubsection*{Loss to unresolved sinks.}

Unresolved sinks are represented by first-order losses:
\begin{equation}
    X_a \longrightarrow \emptyset,
    \qquad
    J_a^{\mathrm{sink}} = D_a c_a .
    \label{eq:sink_loss_general}
\end{equation}
When sinks are unbiased, their contributions cancel in the signed-defect
balance. When sinks are biased, their net contribution must be accumulated as a
separate sink-flux correction.

%===================================================================
\subsection{Operator Form}
\label{sec:operator_form}
%===================================================================

The deterministic master equation may be decomposed into production, binary
collision, unary transition, gas exchange, population-transfer, and sink terms:
\begin{equation}
    \dot{c}_a
    =
    G_a
    +
    \mathcal{Q}_a(\mathbf{c},\mathbf{c})
    +
    \mathcal{L}_{ab}c_b
    +
    \mathcal{H}_a(\mathbf{c},\mathbf{c}^g)
    +
    \mathcal{P}_{ab}c_b
    -
    D_a c_a .
    \label{eq:deterministic_operator_decomposition}
\end{equation}
Here $\mathcal{Q}_a$ collects binary cluster--cluster collisions,
$\mathcal{L}_{ab}$ collects unary emission, trap mutation, and re-solution
operators, $\mathcal{H}_a$ collects gas-reservoir exchange terms, and
$\mathcal{P}_{ab}$ collects population-transfer processes. This decomposition
is only organizational; all of these terms are generated by the single
stoichiometric form in Eq.~\eqref{eq:abstract_master_equation}.

The binary collision operator may be written as
\begin{equation}
    \mathcal{Q}_a(\mathbf{c},\mathbf{c})
    =
    \frac{1}{2}C^+_{abc}K_{bc}c_b c_c
    -
    C^-_{abc}K_{bc}c_b c_c,
    \label{eq:binary_collision_operator}
\end{equation}
where repeated state indices are summed. The tensor $C^+_{abc}$ identifies
binary reactions in which states $b$ and $c$ produce state $a$, while
$C^-_{abc}$ identifies binary reactions in which state $a$ is consumed by
collision with another state. The factor $1/2$ avoids double counting of
symmetric collisions.

%===================================================================
\subsection{Conservation Laws}
\label{sec:general_conservation_laws}
%===================================================================

The abstraction is constrained by two primary conservation laws: signed
lattice-defect balance and gas-species balance.

\subsubsection{Signed lattice-defect balance.}

Define the signed lattice-defect inventory
\begin{equation}
    Q(t)
    =
    \sum_{n,\alpha} n c^I_{n,\alpha}(t)
    -
    \sum_{m,\boldsymbol{\ell},\beta}
    m c^V_{m,\boldsymbol{\ell},\beta}(t)
    +
    Q_{\mathrm{sink}}(t) .
    \label{eq:signed_defect_inventory}
\end{equation}
The cumulative term $Q_{\mathrm{sink}}$ accounts for biased absorption by
unresolved sinks. Differentiating gives
\begin{equation}
    \frac{dQ}{dt}
    =
    \sum_{n,\alpha} n G^I_{n,\alpha}
    -
    \sum_{m,\boldsymbol{\ell},\beta} m G^V_{m,\boldsymbol{\ell},\beta}
    +
    \dot{Q}_{\mathrm{bias}} .
    \label{eq:signed_defect_balance}
\end{equation}
For conservative Frenkel-pair production and unbiased sinks,
\begin{equation}
    \frac{dQ}{dt}=0 .
\end{equation}
When biased sink absorption is present,
\begin{equation}
    Q_{\mathrm{sink}}(t)
    =
    \int_0^t
    \left[
    J_I^{\mathrm{bias}}(t')
    -
    J_V^{\mathrm{bias}}(t')
    \right]dt' .
    \label{eq:biased_sink_correction}
\end{equation}
This is the signed-defect version of the swelling balance: vacancy storage in
cavities is balanced by interstitial storage in loops and the cumulative net
biased absorption of interstitials by microstructural sinks.

\subsubsection{Gas-species balance.}

For each gas species $\mu$, define the total gas inventory
\begin{equation}
    M_\mu(t)
    =
    c^g_\mu(t)
    +
    \sum_{m,\boldsymbol{\ell},\beta}
    \ell_\mu c^V_{m,\boldsymbol{\ell},\beta}(t).
    \label{eq:gas_inventory_general}
\end{equation}
Gas trapping, gas emission, re-solution, vacancy-cluster coalescence, trap
mutation, population transfer, and signed-defect cancellation redistribute gas
between the matrix reservoir and vacancy clusters, but do not create or destroy
gas. Therefore,
\begin{equation}
    \boxed{
    \frac{dM_\mu}{dt}
    =
    G^g_\mu
    -
    J^{g,\mathrm{ext}}_\mu .
    }
    \label{eq:gas_species_conservation}
\end{equation}
Here $G^g_\mu$ is the external production rate of gas species $\mu$, and
$J^{g,\mathrm{ext}}_\mu$ is an explicitly modeled external gas loss. In a
closed matrix reservoir,
\begin{equation}
    J^{g,\mathrm{ext}}_\mu=0,
    \qquad
    \frac{dM_\mu}{dt}=G^g_\mu .
\end{equation}

%===================================================================
\subsection{Stochastic Jump-Process Counterpart}
\label{sec:stochastic_jump_process_counterpart}
%===================================================================

The same reaction-network structure defines a stochastic cluster-dynamics
model. Let $N_a(t)$ denote the integer number of objects in state $X_a$ in a
representative volume $\Omega$. Each reaction $r$ changes the stochastic state
by
\begin{equation}
    \mathbf{N}
    \longrightarrow
    \mathbf{N}+\mathbf{S}_r,
\end{equation}
where $\mathbf{S}_r$ is the stoichiometric increment vector for reaction $r$.
The stochastic propensity of reaction $r$ is denoted by $a_r(\mathbf{N})$.
For a unary reaction,
\begin{equation}
    a_r(\mathbf{N}) = \Gamma_r N_b .
\end{equation}
For a binary reaction between distinct states $b$ and $c$,
\begin{equation}
    a_r(\mathbf{N})
    =
    \frac{K^{(r)}_{bc}}{\Omega}N_bN_c .
\end{equation}
For a self-collision $b=c$,
\begin{equation}
    a_r(\mathbf{N})
    =
    \frac{K^{(r)}_{bb}}{\Omega}
    \frac{N_b(N_b-1)}{2} .
\end{equation}
The deterministic concentration equation is recovered in the mean-field limit
\begin{equation}
    c_a = \frac{N_a}{\Omega},
    \qquad
    \Omega\rightarrow\infty .
\end{equation}
Thus the deterministic and stochastic formulations use the same state list,
reaction list, stoichiometric matrix, and conservation checks.

%===================================================================
\subsection{Code-Generation-Ready Representation}
\label{sec:code_generation_ready_representation}
%===================================================================

For implementation, each state can be represented by a structured record
\begin{equation}
    a
    =
    \left(
    \mathrm{kind},
    q,
    n,
    m,
    \boldsymbol{\ell},
    p,
    \mathrm{mobility},
    \mathrm{sink\_coupling}
    \right),
    \label{eq:state_record}
\end{equation}
where
\begin{equation}
    \mathrm{kind}
    \in
    \{\mathrm{I\_cluster},\mathrm{V\_cavity},\mathrm{gas\_reservoir}\},
\end{equation}
$q$ is the signed lattice-defect content, $n$ is the SIA number if the state is
an interstitial cluster, $m$ is the vacancy number if the state is a vacancy
cluster, $\boldsymbol{\ell}$ is the gas occupancy vector, and $p$ is the
population label. For interstitial states, $p=\alpha$; for vacancy states,
$p=\beta$; for gas-reservoir states, $p$ may be omitted.

Each reaction can be represented by
\begin{equation}
    r
    =
    \left(
    \mathcal{R}^-_r,
    \mathcal{R}^+_r,
    J_r,
    \mathbf{S}_r,
    \mathcal{C}_r
    \right),
    \label{eq:reaction_record}
\end{equation}
where $\mathcal{R}^-_r$ is the reactant list, $\mathcal{R}^+_r$ is the product
list, $J_r$ is the rate law, $\mathbf{S}_r$ is the stoichiometric increment,
and $\mathcal{C}_r$ is a set of conservation checks.

For an internally conservative lattice-defect reaction, the signed-defect
check is
\begin{equation}
    \sum_a q_a S_{ar}=0 .
    \label{eq:code_signed_defect_check}
\end{equation}
For an internally conservative gas-transfer reaction, the gas-species checks
are
\begin{equation}
    \sum_a \ell_{\mu a} S_{ar}=0,
    \qquad
    \mu=1,\ldots,N_g .
    \label{eq:code_gas_check}
\end{equation}
External sources and sinks are not required to satisfy the zero-sum checks, but
their contributions must enter the corresponding production, loss, or
cumulative reservoir terms.

The deterministic ODE system is assembled as
\begin{equation}
    \dot{\mathbf{c}}
    =
    \mathbf{G}
    +
    \mathbf{S}\mathbf{J}(\mathbf{c})
    -
    \mathbf{D}\mathbf{c},
    \label{eq:assembled_ode_system}
\end{equation}
where $\mathbf{S}$ is sparse. The same data structure defines stochastic
propensities for stochastic cluster dynamics.

%===================================================================
\subsection{Specialization to the Vacancy--Interstitial--Helium System}
\label{sec:specialization_V_I_He}
%===================================================================

The vacancy--interstitial--helium system is recovered by setting
\begin{equation}
    N_g=1,
    \qquad
    g_1=\mathrm{He},
    \qquad
    \boldsymbol{\ell}=(\ell).
\end{equation}
If only one interstitial-cluster population and one vacancy-cluster population
are retained, then
\begin{equation}
    N_P^I=1,
    \qquad
    N_P^V=1,
\end{equation}
and the population labels can be suppressed. The general variables reduce to
\begin{equation}
    c^I_{n,\alpha}\rightarrow c_n,
    \qquad
    c^V_{m,(\ell),\beta}\rightarrow c_{m,\ell},
    \qquad
    c^g_1\rightarrow c_h .
    \label{eq:VIHe_variable_reduction}
\end{equation}
Thus the two resolved cluster populations are
\begin{equation}
    \{I_n\}_{n=1}^{N_I},
    \qquad
    \{V_m\mathrm{He}_{\ell}\}_{m,\ell},
\end{equation}
and the matrix gas reservoir is the free helium concentration $c_h$.

Gas trapping becomes
\begin{equation}
    \mathrm{He}
    +
    V_m\mathrm{He}_{\ell}
    \longrightarrow
    V_m\mathrm{He}_{\ell+1},
\end{equation}
while helium emission and radiation re-solution become
\begin{equation}
    V_m\mathrm{He}_{\ell}
    \longrightarrow
    V_m\mathrm{He}_{\ell-1}
    +
    \mathrm{He} .
\end{equation}
The signed interstitial--vacancy collision rule becomes
\begin{equation}
    I_n + V_m\mathrm{He}_{\ell}
    \longrightarrow
    \begin{cases}
    I_{n-m}+\ell\,\mathrm{He}, & n>m, \\[4pt]
    \ell\,\mathrm{He}, & n=m, \\[4pt]
    V_{m-n}\mathrm{He}_{\ell}, & m>n .
    \end{cases}
    \label{eq:VIHe_signed_collision}
\end{equation}
If the vacancy cluster survives, its helium content remains trapped. If the
cavity is eliminated, the helium inventory is returned to the matrix reservoir.

Trap mutation becomes
\begin{equation}
    V_m\mathrm{He}_{\ell}
    \longrightarrow
    V_{m+1}\mathrm{He}_{\ell}
    +
    I_1,
    \label{eq:VIHe_trap_mutation}
\end{equation}
which conserves signed lattice-defect content because
\begin{equation}
    -m=-(m+1)+1 .
\end{equation}

The helium conservation law reduces to
\begin{equation}
    \frac{d}{dt}
    \left[
    c_h
    +
    \sum_{m=1}^{N_V}
    \sum_{\ell=1}^{\ell_{\max}(m)}
    \ell c_{m,\ell}
    \right]
    =
    G_{\mathrm{He}}
    -
    J_{\mathrm{He}}^{\mathrm{ext}} .
    \label{eq:He_conservation_specialized}
\end{equation}
For a closed matrix reservoir with no external helium escape,
\begin{equation}
    J_{\mathrm{He}}^{\mathrm{ext}}=0 .
\end{equation}

The signed Frenkel-pair balance reduces to
\begin{equation}
    Q(t)
    =
    \sum_{n=1}^{N_I} n c_n(t)
    -
    \sum_{m=1}^{N_V}
    \sum_{\ell=0}^{\ell_{\max}(m)}
    m c_{m,\ell}(t)
    +
    Q_{\mathrm{sink}}(t) .
    \label{eq:VIHe_signed_balance}
\end{equation}
For conservative cascade production and unbiased sinks, $Q(t)$ is constant.
With biased dislocation absorption, $Q_{\mathrm{sink}}(t)$ is the cumulative
net biased absorption term. This is the signed-defect form of the swelling
identity.
